\section{Introducción}
La transferencia de calor y el establecimiento del equilibrio térmico permiten describir la evolución de cuerpos que se encuentran inicialmente a distintas temperaturas. Estos procesos son fundamentales en termodinámica porque relacionan los cambios de estado de los sistemas con los intercambios de energía. La calorimetría proporciona un método experimental para estudiar dichos intercambios mediante mediciones de masa y temperatura. Cuando el intercambio con el entorno y el trabajo son despreciables, la conservación de energía relaciona el calor cedido por unos componentes con el absorbido por los demás \cite{serway}.

En una experiencia de mezcla, el recipiente también participa en el balance térmico. Considerar la capacidad calorífica del calorímetro permite representar su contribución y evaluar las limitaciones de un modelo que atribuye todo el intercambio a las sustancias contenidas. Cuando se incorpora hielo, el balance debe incluir tanto la energía necesaria para la fusión como el calentamiento posterior del agua resultante. El calor latente permite describir esa transferencia de energía durante el cambio de fase, que no se manifiesta como una variación de temperatura mientras ocurre la fusión en equilibrio \cite{serway}.

La presente práctica estudió estos procesos en dos etapas: la estimación de la capacidad calorífica del calorímetro mediante mezclas de agua y la comparación de temperaturas de equilibrio en un sistema agua-hielo con modelos que consideran o desprecian dicha capacidad. El tratamiento de las incertidumbres permitió examinar el alcance de las mediciones y las limitaciones de la comparación entre teoría y experimento \cite{taylor}. El trabajo se orientó a los objetivos de la guía de laboratorio \cite{guiauca4}.

\subsection{Objetivos}
\textbf{Objetivo general.} Analizar experimentalmente el intercambio de energía térmica y el establecimiento del equilibrio térmico mediante experiencias de calorimetría, considerando la capacidad calorífica del calorímetro y el cambio de fase del agua.

\textbf{Objetivos específicos.}
\begin{enumerate}
  \item Determinar experimentalmente la capacidad calorífica del calorímetro mediante el balance energético de mezclas de agua a diferentes temperaturas.
  \item Comparar la temperatura de equilibrio experimental del sistema agua-hielo con las predicciones teóricas que consideran y desprecian la capacidad calorífica del calorímetro.
  \item Analizar las principales fuentes de error e incertidumbre que limitan las mediciones y la interpretación de los resultados de calorimetría.
\end{enumerate}
